\chapter*{Preface}
\addcontentsline{toc}{chapter}{Preface}
\markboth{Preface}{Preface}

Take a positive integer. If it is even, halve it; if it is odd, triple it and
add one. Repeat. The Collatz conjecture says that every starting value
eventually reaches $1$. It can be explained to a child in a minute, checked by
a laptop for every number a person could write down in a lifetime, and it has
resisted every serious attempt at proof since the 1930s. This book is about
why.

We came to the problem from outside number theory. One of us works in Bayesian
statistics and probability, the other in optimization and machine learning,
and what drew us in was a simple observation that anyone trained in either
field makes within an hour of looking at the map. The logarithm of the iterate
behaves like a random walk with negative drift. Each odd step multiplies by
about $3/2$, each even step divides by $2$, and if the parities were fair coin
flips the expected change in $\log x$ per step would be
$\tfrac12\log\tfrac32-\tfrac12\log2=\log(\sqrt3/2)\approx-0.144$. That is a
Lyapunov function decreasing in expectation, the most familiar object in
optimization and control, and it predicts the growth of the stopping times of
real orbits to within about two percent. None of it is a proof, and the reason it is not a
proof is the subject of the book.

\section*{What the book does}

We read the Collatz dynamics as an optimization problem and push the reading
as far as it will go. Three things come out.

The first is an accounting of what a proof would have to contain. The
conjecture splits into two halves that require completely different tools.
The divergence half asks whether $\log x_n$ can drift upward forever along a
single orbit.
It is equivalent to a statement about Ces\`aro averages of the dyadic valuations $\vt(3y+1)$ along that orbit, namely that the lower limit of their running mean exceeds $\log_2 3\approx1.585$ (\cref{conj:cesaro}), and it is a one sided law of large numbers for a deterministic sequence that nobody knows how to prove.
The cycle half asks
whether an orbit can close up. It is not an averaging question at all: a cycle
forces a product of numbers slightly larger than $3$ to equal a power of $2$,
so the rational approximations of $\log_2 3$ enter, although, as we show, they cannot decide it on their own.

The second is a set of structural theorems, most of them short and several of
them classical. We prove that the Syracuse
exponent is exactly Geometric$(\tfrac12)$ on residue classes, that under the
Bernoulli model the iterates form a multiplicative martingale while their
logarithms form a strict supermartingale (the gap between the two is Jensen's
inequality), that the parity autocorrelation function of the standard map is
$\rho(h)=(-\tfrac12)^h$ in the model, and that no monotone Lyapunov function certifies descent over any
fixed number of strides. On cycles we give a short mediant lemma for the one sided distances $\lceil r\log_23\rceil-r\log_23$ which, with a self contained proof of the classical description of the one sided best approximations, shows that only the upper rungs of the semiconvergent ladder of $\log_23$ ever need to be checked. Combined with verification up to $2^{68}$ it gives, unconditionally,
that every nontrivial cycle has at least $72\,057\,431\,991$ odd elements. The
method goes back to Crandall and Eliahou, the bound is that of Hercher and
Puchert, and we claim no priority for any numerical bound; the current record,
$r\geq1.375\times10^{11}$, is Hercher's, made unconditional by Ba\v{r}ina's
verification to $2^{71}$. What we offer is a short rigorous route to the bound
and an exact integer certificate for it. We then prove that Diophantine
information about $\log_23$ alone cannot finish the cycle half, since infinitely many cycle lengths survive
every verification range, and we show with a window lemma that an arithmetic
improvement of the basic cycle inequality by a factor slightly larger than $r/3$, combined with a finite verification, would.
On the computational side, nested, split and vertical likelihood sampling reach regions of exponentially small prior mass in a number of steps polynomial in the depth, provided each level set can be sampled exactly. We prove that a perfect version of split sampling samples the exceptional residue classes exactly, at any depth and in time polynomial in the depth. Beyond the bit length of the integers the level set samplers eventually stop, and for split sampling the census shows why: in the boxes we computed, the slow integers are a small number of long orbits together with their tributaries, and a local sampler leaves the tree it starts in only about as often as rejection sampling finds a new one.

The third is an algorithm. The Collatz map alternates an ascent, the odd step,
with a corrective descent, the run of halvings, and the ratio of the two
logarithmic step sizes is $\beta=\log(3/2)/\log2=\log_23-1\approx0.585$. We
turn this into Switching MM, a first order method with fixed momentum $\beta$
that extrapolates only on odd iterations. On quadratics we prove that two of
its steps are exactly a gradient step of size $\alpha$ followed by one of size
$(1+\beta)\alpha$, so that it is a cyclic step size method in disguise, and we
compute its stationary variance under gradient noise in closed form. The
noise amplification of Switching MM relative to gradient descent at the same step is bounded by $5/3$ for every $\beta<1$,
whereas in the slow modes every step momentum amplifies noise by $1/(1-\beta)$, which is $10$ at
the machine learning default $\beta=0.9$. That is the precise sense in which
the parity schedule is a low tuning default. For any smooth strongly convex
function we prove a linear rate whose factor is the exact worst case for the distance to the minimizer, and a gain weight theorem shows that no periodic
momentum schedule of any period can beat gradient descent's noise floor in the slow modes at
matched speed: momentum trades speed for noise in exact proportion, and unequal
weighting of gradients pays a Cauchy and Schwarz penalty.

\section*{What the book does not do}

It does not prove the Collatz conjecture. We say this at the outset and repeat
it where it matters, because the literature on the problem contains many
claimed proofs and the optimization language makes it easy to mistake a
heuristic for an argument. Several of the most instructive passages in the
book are pitfalls: arguments that look like progress and are not. The Ces\`aro
criterion $\bar k_N>\log_23$ holds automatically at the stride $N$ at which an orbit reaches $1$,
so checking it numerically on terminating orbits carries no information. Convergent
denominators alone do not locate the records of the cycle bound, because the
relevant distance is one sided. In the stochastic model, the linear majorizer of the logarithm has positive mean increment while the logarithm itself has negative drift.
Two constants near one half, a momentum rate that pairs an ascent and a descent measured at different levels and the modulus of an eigenvalue of a chain conditioned on a set that it leaves, agree by accident (\cref{rem:ch5:mixing}).

\section*{How to read it}

Part I sets up the problem: the history and the three standard maps (\cref{chap:intro}), the arithmetic of parity vectors and the Syracuse map including Terras' theorem and the density result (\cref{chap:parity}), and the stochastic models with their martingales (\cref{chap:stochastic}). Part II is the optimization reading: Lyapunov functions and their limits (\cref{chap:lyapunov}); majorization and momentum (\cref{chap:momentum}); majorization pressed to its limit on the Collatz map, where it reproduces the cycle inequality and shows that no correction of $\log x$ by a function of the last $d$ binary digits of $x$, for any fixed $d$, certifies descent (\cref{chap:mmcollatz}); Switching MM as an optimizer (\cref{chap:smm}) and its behaviour under noise, with experiments (\cref{chap:noise}); nonlinear difference equations and momentum MM, with difference inequalities for the inverse tree and a linear rate for Switching MM on smooth strongly convex functions that is tight for the distance to the minimizer (\cref{chap:difference}); and the three point inequalities and similar triangles behind every convergence proof, which give Switching MM a convex rate a factor $1+\beta/2$ better than gradient descent and show what the Collatz map lacks (\cref{chap:threepoint}). Part III treats cycles through continued fractions, with a section on the sibling problem of Erd\H{o}s $125$ settled by machine checked proof search and on our own Lean statements (\cref{chap:cycles}); the divergence half through Ces\`aro averages, mixing and Tao's theorem (\cref{chap:divergence}); the positive density theorems of 2026, read as the stride model run backwards through the inverse tree, with numerical checks of their constants and of the $3$ adic law they rest on (\cref{chap:backward}); a synthesis of what a proof would require (\cref{chap:path}); a computational path, with what the MM view implies for automated certificates (\cref{chap:computation}); what nested, split and vertical likelihood sampling can and cannot do, with the tree structure of the slow integers (\cref{chap:samplers}); progress on six of the open problems, including the third cumulant of stopping times and the exact convex rate of Switching MM (\cref{chap:progress}); a summary of what the book contributes, a ledger of the strategies it rules out, and a computational road map to a solution (\cref{chap:synthesis}); and a list of open problems and computational projects (\cref{chap:open}). A gallery of computed pictures follows Part III, and two appendices tabulate the constants and describe the code.

A reader interested only in the optimizer can read \cref{chap:momentum,chap:smm,chap:noise} and then \cref{chap:difference,chap:threepoint}, which use their results, turning back to Part I and \cref{chap:lyapunov} for the stride arithmetic behind the momentum constant. A reader interested only in the number theory can go from Part I directly to \cref{chap:cycles,chap:divergence}, but before \cref{chap:path,chap:computation} needs the impossibility results of \cref{chap:lyapunov,chap:mmcollatz} and the first two sections of \cref{chap:difference} (with \cref{sec:ch5b:stability} for two comparisons in \cref{chap:computation}). Every chapter ends with notes, and every chapter except \cref{chap:open} with exercises.

\section*{Acknowledgements}

We thank participants at the Chicago Booth Statistics Workshop and colleagues
who listened to early versions of these arguments and pointed out the ones
that were wrong. The computations were done in Python with mpmath, NumPy,
SciPy, SymPy, Numba, CVXPY and scikit-learn, and the figures were drawn with Matplotlib; the formal statements in \cref{chap:cycles} were checked in
Lean~4.

\vspace{12pt}
\noindent Chicago and Fairfax, October 2026\hfill N.\,G.\,P.\quad V.\,S.
